\cotefile{57}
\scanpage{133}
% lines 1–5
\uline{Théorème}$^{\ast}$ Soit \deleted{$\unsure{Y'\!-\!X'|Y'\leftarrow X'|X'\to X''}$} \inserted{$f:X\unsure{-|\to}Y$} un morphisme \uline{surjectif} de préschémas propres sur un corps $k$. Alors le noyau $\underline{\underline{N}}$ de $\underline{\underline{\mathrm{Pic}}}_{\deleted{X}Y/k}-\underline{\underline{\mathrm{Pic}}}_{X/k}$ est \margin{?}\deleted{contenu} un sous-groupe d'un \uline{groupe} \uline{affine} $G$ sur $k$.

% lines 6–8
\unsure{Donc}

\uline{Corollaire} Si $\underline{\mathrm{Pic}}_{Y/k}$ et $\underline{\mathrm{Pic}}_{X/k}$ existent, le noyau de $\underline{\mathrm{Pic}}_{Y/k}-\underline{\mathrm{Pic}}_{X/k}$ est un groupe alg. affine.

% lines 9–14
\uline{Corollaire} Sous les conditions précédentes, \unsure{on} $\text{Nér Sév}_{0}(Y/\unsure{k})-\text{Nér Sév}_{0}(X/k)$ est \uline{injectif}, et \deleted{\unsure{Q}} si $A$ et $B$ sont les ``parties abéliennes'' de $\underline{\mathrm{Pic}}_{Y/k}$ et $\underline{\mathrm{Pic}}_{X/k}$ [définies si $k$ est parfait] alors $A\to B$ a un noyau fini.

% lines 15–19
\uline{Dém} \deleted{$\underline{\mathrm{Pic}}_{X}-\underline{\mathrm{Pic}}_{Y}$} On va pour simplifier admettre que tous les $\underline{\mathrm{Pic}}$ envisagés sont représentables. \inserted{a) On peut supposer $k$ \uline{alg. clos} (\unsure{Sert à rien|Suffira|ou non} ?)} b) On aura un diagramme de \uuline{schémas} sur $k$
\diagram{\begin{tikzcd}[ampersand replacement=\&]
\underline{\mathrm{Pic}}_{X/k}\arrow[d,dash] \& \underline{\mathrm{Pic}}_{Y/\unsure{k}}\arrow[l]\arrow[d,"\text{\uline{affine}}"] \& {} \\
\underline{\mathrm{Pic}}_{X_{\mathrm{réd}}} \& \underline{\mathrm{Pic}}_{Y_{\mathrm{red}}}\arrow[l,"?"'] \& {}
\end{tikzcd}\(\left.\vphantom{\begin{matrix}X\\X\\X\end{matrix}}\right]\text{ par \unsure{vert|vert.|dont}}\)}
% archive sheet number: 132

